% =====================================================================
\chapter{Cuarto instrumento: los regímenes de mercado}
\chaptermark{Los regímenes de mercado}
\label{cap:hmm}

El VIX argentino dice cuánta tensión hay cada día, pero se mueve de forma muy dispareja: hay
tramos largos de calma y episodios cortos y violentos de estrés. Para trabajar con esa
alternancia hace falta una clasificación de los días según el estado del mercado. El cuarto
instrumento la produce con un modelo oculto de Markov: un modelo que supone que el mercado
pasa por unos pocos estados que no se observan directamente, y que los infiere de la
volatilidad.

% ---------------------------------------------------------------------
\section{El modelo}

Un modelo de cambio de régimen de Markov supone que la serie —aquí, el VIX argentino— tiene una
media y una varianza distintas en cada uno de $k$ estados ocultos, y que el mercado pasa de un
estado a otro según una matriz de probabilidades de transición
\citep{hamilton1989,kim1994}: la probabilidad de estar en un régimen mañana depende solo del
régimen de hoy. En la literatura de mercados se lo suele llamar HMM \citep{rabiner1989}. Se
estimó con \texttt{MarkovRegression} de \texttt{statsmodels}, con media y varianza propias de
cada régimen, sobre las 323 ruedas del VIX argentino.

A cada día se le asigna el régimen de mayor probabilidad \textit{suavizada}, que usa toda la
serie. Se guarda también la versión \textit{filtrada}, que solo usa la información disponible
hasta ese día y es la adecuada si el régimen se usa para pronosticar \citep{majumder2023}. Las
dos coinciden en el 98,5\% de los días.

% ---------------------------------------------------------------------
\section{¿Cuántos regímenes?}

El número de regímenes no se fijó a mano: se ajustó el modelo para $k=2$ a 6 y se comparó el
criterio bayesiano de información \citep[BIC;][]{schwarz1978}, que premia el ajuste y penaliza
los parámetros de más. Hubo que agregar una regla: descartar los ajustes que no convergen o que
crean un régimen de menos de cinco días. Un régimen de un solo día tiene varianza casi cero, su
verosimilitud se dispara y el BIC mejora de forma artificial: es una falla conocida de los
modelos de mezcla, no un estado real del mercado (tabla~\ref{tab:bic}).

\begin{table}[H]
\centering\small
\begin{tabularx}{\linewidth}{@{}cccccY@{}}
\toprule
$k$ & \textbf{BIC} & \textbf{Convergió} & \textbf{Régimen más chico} & \textbf{Válido} & \textbf{Qué pasa} \\
\midrule
2 & 2.202,9 & No & 1 día & No & Depende del arranque aleatorio \\
\textbf{3} & \textbf{2.119,4} & \textbf{Sí} & \textbf{65 días} & \textbf{Sí} & \textbf{Elegido; idéntico con 10 semillas} \\
4 & 2.152,0 & No & 1 día & No & Los mismos 3 regímenes más un día suelto (01/12/2025) \\
5 & 2.161,4 & No & 1 día & No & Un estado clavado en el mínimo de la serie \\
6 & 2.185,2 & No & 1 día & No & Ídem \\
\bottomrule
\end{tabularx}
\caption{Selección del número de regímenes. Fuente: elaboración propia
(\carpeta{Lado Bursatil/06_hmm_regimen_mercado}).}
\label{tab:bic}
\end{table}

Con $k=3$ el BIC es el menor entre los ajustes válidos, y el resultado es el mismo con diez
semillas distintas: no depende del punto de partida (figura~\ref{fig:bic}).

\begin{figure}[H]
\centering
\includegraphics[width=0.5\linewidth]{bic_curva.png}
\caption{BIC según la cantidad de regímenes. Fuente: elaboración propia
(\carpeta{Lado Bursatil/06_hmm_regimen_mercado}).}
\label{fig:bic}
\end{figure}

% ---------------------------------------------------------------------
\section{Los tres regímenes}

La tabla~\ref{tab:regimenes} describe los tres regímenes, que se nombraron por su nivel de
volatilidad.

\begin{table}[H]
\centering\small
\begin{tabular}{@{}lrrrr@{}}
\toprule
\textbf{Régimen} & \textbf{VIX promedio} & \textbf{Retorno diario promedio} & \textbf{Días} & \textbf{\% del período} \\
\midrule
Calmo & 33,0\% & $-$0,074\% & 177 & 54,8 \\
Moderado & 46,6\% & 0,053\% & 81 & 25,1 \\
Tenso & 70,8\% & 0,518\% & 65 & 20,1 \\
\bottomrule
\end{tabular}
\caption{Regímenes del mercado entre el 2/1/2025 y el 30/4/2026 (323 ruedas). Fuente:
elaboración propia (\carpeta{Lado Bursatil/06_hmm_regimen_mercado}).}
\label{tab:regimenes}
\end{table}

Los regímenes son muy persistentes: la probabilidad de seguir en el mismo régimen al día
siguiente es de 98,0\% en Calmo, 95,8\% en Moderado y 95,1\% en Tenso
(tabla~\ref{tab:transicion}). Más interesante es cómo se pasa de uno a otro. La salida del
estrés es gradual: de Tenso nunca se volvió directo a Calmo, sino siempre a través de Moderado.
La entrada, en cambio, no: las dos veces que el mercado salió de Calmo saltó directo a Tenso,
y las dos fueron el día hábil siguiente a una elección (08/09 y 27/10/2025). El régimen Tenso
tiene el retorno promedio más alto porque incluye la fuerte suba posterior a las legislativas:
la tensión no es lo mismo que la caída.

\begin{table}[H]
\centering\small
\begin{tabular}{@{}lccc@{}}
\toprule
\textbf{De \flecha{} a} & \textbf{Calmo} & \textbf{Moderado} & \textbf{Tenso} \\
\midrule
Calmo & 98,0\% & 0,0\% & 2,0\% \\
Moderado & 4,2\% & 95,8\% & 0,0\% \\
Tenso & 0,0\% & 4,9\% & 95,1\% \\
\bottomrule
\end{tabular}
\caption{Probabilidad de pasar de un régimen (fila) a otro (columna) de una rueda a la
siguiente. Fuente: elaboración propia (\carpeta{Lado Bursatil/06_hmm_regimen_mercado}).}
\label{tab:transicion}
\end{table}

La figura~\ref{fig:hmm} muestra la secuencia completa. El período empieza en Moderado, pasa a
Tenso el 07/04/2025 con los aranceles y la salida del cepo, y se calma desde fines de mayo. Los
dos episodios electorales son tramos Tenso de exactamente 20 ruedas, del 08/09 al 03/10 y del
27/10 al 21/11: el modelo hereda el eco del VIX (capítulo~\ref{cap:vix}). Desde el 01/12/2025
hasta el final de la ventana el mercado queda en Calmo, durante 100 ruedas seguidas.

\begin{figure}[H]
\centering
\includegraphics[width=0.8\linewidth]{hmm_regimen_mercado.png}
\caption{S\&P~Merval y VIX argentino coloreados por régimen, y probabilidad suavizada de cada
régimen por día. Fuente: elaboración propia (\carpeta{Lado Bursatil/06_hmm_regimen_mercado}).}
\label{fig:hmm}
\end{figure}

\textbf{Para qué sirve en la tesina.} El régimen es la variable de contexto con la que se
pregunta si la relación entre noticias y precios cambia según el estado del mercado
\citep{garcia2013}. En la ventana común con las noticias (244 ruedas) hay 177 ruedas Calmo, 18
Moderado y 49 Tenso; como Moderado queda con muy pocas ruedas, el análisis por régimen del
capítulo~\ref{cap:var} compara Calmo contra Tenso.
